\documentclass[11pt,a4paper,oneside,openany]{book}
\usepackage[utf8]{inputenc}
\usepackage[english]{babel}
\usepackage[margin=1in]{geometry}
\usepackage{graphicx}
\usepackage{tikz}
\usepackage{pgfplots}
\usepackage{xcolor}
\usepackage{fancyhdr}
\usepackage{listings}
\usepackage{hyperref}
\usepackage{array}
\usepackage{booktabs}
\usepackage{amsmath}
\usepackage{amssymb}
\usepackage{float}
\usepackage{caption}
\usepackage{subcaption}
\usepackage{longtable}
\usepackage{multirow}
\usepackage{setspace}
\usepackage{tcolorbox}
\usepackage{mdframed}
\usepackage{soul}

% TikZ libraries
\usetikzlibrary{shapes,arrows,positioning,fit,calc,shadows,patterns,decorations.pathreplacing}

% Color scheme
\definecolor{darkblue}{RGB}{26,84,144}
\definecolor{lightblue}{RGB}{45,122,184}
\definecolor{accent}{RGB}{74,144,200}
\definecolor{background}{RGB}{240,245,250}
\definecolor{codebackground}{RGB}{245,245,245}

% Header and footer
\pagestyle{fancy}
\fancyhf{}
\fancyhead[R]{\textcolor{darkblue}{\slshape HighPerformance}}
\fancyfoot[C]{\thepage}
\renewcommand{\headrulewidth}{0.5pt}
\renewcommand{\headrule}{\color{darkblue}\hrule}
\renewcommand{\footrulewidth}{0.5pt}

% Code listing setup
\lstset{
    basicstyle=\ttfamily\small,
    breaklines=true,
    showstringspaces=false,
    tabsize=2,
    language=JavaScript,
    backgroundcolor=\color{codebackground},
    keywordstyle=\color{darkblue}\bfseries,
    commentstyle=\color{gray}\itshape,
    stringstyle=\color{accent},
    breakatwhitespace=true,
    frame=single,
    frameround=tttt,
    rulecolor=\color{lightblue},
    captionpos=b,
    float=h!,
    aboveskip=10pt,
    belowskip=10pt
}

% Section formatting
\usepackage{titlesec}
\titleformat{\chapter}[block]{\Large\bfseries\color{darkblue}}{Chapter \thechapter}{1em}{}[\vspace{0.5pt}\color{darkblue}\rule{\textwidth}{1.5pt}]
\titleformat{\section}{\large\bfseries\color{lightblue}}{\thesection}{1em}{}
\titleformat{\subsection}{\bfseries\color{accent}}{\thesubsection}{1em}{}

% Custom commands
\newcommand{\todo}[1]{\textcolor{red}{TODO: #1}}
\newcommand{\note}[1]{\begin{tcolorbox}[colback=blue!5!white,colframe=darkblue,title=Note]\small #1\end{tcolorbox}}
\newcommand{\warning}[1]{\begin{tcolorbox}[colback=orange!5!white,colframe=orange,title=Warning]\small #1\end{tcolorbox}}
\newcommand{\code}[1]{\texttt{\small #1}}
\newcommand{\endpoint}[1]{\fcolorbox{lightblue}{background}{\texttt{#1}}}

% Title page
\title{\Huge\textbf{\color{darkblue}HighPerformance}\\[0.5cm]
        \Large Performance Management \& Bonus Calculation System\\[2cm]
        \includegraphics[width=0.3\textwidth]{example-image-a}\\[2cm]
        \normalsize Project Documentation\\
        \textit{Version 1.0 --- February 2026}}
\author{\textbf{Development Team}\\Integration Architecture Project\\[1cm]\textcolor{darkblue}{\Large Enterprise-Grade System}}
\date{\today}

\begin{document}

% Title page
\maketitle

% Abstract
\chapter*{Abstract}
\addcontentsline{toc}{chapter}{Abstract}

HighPerformance is a comprehensive enterprise-grade system designed for sales organizations to automate bonus calculations, manage performance metrics, and provide real-time dashboards with multi-level approval workflows. This document provides complete technical documentation including architecture, API reference, database schemas, integration guides, and deployment procedures.

\begin{center}
\begin{tabular}{ll}
\toprule
\textbf{Property} & \textbf{Value} \\
\midrule
Project Type & Web Application \\
Technology Stack & Node.js, Angular, MongoDB, Camunda \\
Status & Active Development \\
Last Updated & February 11, 2026 \\
\bottomrule
\end{tabular}
\end{center}

% Table of contents
\tableofcontents
\newpage

% ===== CHAPTER 1: EXECUTIVE SUMMARY =====
\chapter{Executive Summary}

\section{Project Overview}

HighPerformance addresses enterprise requirements for performance management and bonus automation. The system provides:

\begin{itemize}
    \item \textbf{Automated Bonus Calculation} --- Multi-component computations based on performance metrics
    \item \textbf{Multi-Level Approval Workflow} --- CEO and HR approval with process orchestration
    \item \textbf{Real-Time Dashboards} --- Role-based analytics and visualization
    \item \textbf{Enterprise Integration} --- Connectors for OrangeHRM, OpenCRX, Odoo
    \item \textbf{Security-First Design} --- JWT authentication, RBAC, encrypted communications
    \item \textbf{Audit & Compliance} --- Complete logging and historical tracking
\end{itemize}

\section{Technology Stack}

\begin{table}[H]
\centering
\begin{tabular}{lll}
\toprule
\textbf{Layer} & \textbf{Technology} & \textbf{Version} \\
\midrule
Frontend & Angular & 16+ \\
Backend & Express.js (Node.js) & 4.x \\
Database & MongoDB & 6.x+ \\
Workflow Engine & Camunda BPM & Latest \\
Authentication & JWT & jsonwebtoken 9.0+ \\
API Documentation & OpenAPI & 3.0 \\
Containerization & Docker & Latest \\
\bottomrule
\end{tabular}
\caption{Technology Stack Components}
\label{tab:techstack}
\end{table}

\section{Key Stakeholders}

\begin{tcolorbox}[colback=background,colframe=darkblue,title=Supported Roles]
\begin{description}
    \item[CEO] Compute bonuses, approve CEO-level, access executive dashboards
    \item[HR] Manage salesmen, create performance records, approve HR-level
    \item[Salesman] View personal bonuses (when released), submit qualifications
    \item[Admin] System configuration, user management, integrations
\end{description}
\end{tcolorbox}

% ===== CHAPTER 2: ARCHITECTURE =====
\chapter{System Architecture}

\section{High-Level Architecture}

\begin{figure}[H]
\centering
\begin{tikzpicture}[scale=0.7, node distance=3.5cm]
    % Define styles
    \tikzstyle{client}=[rectangle,rounded corners,draw=darkblue,fill=background,thick,minimum width=2.2cm,minimum height=1.1cm,text width=2cm,align=center,font=\small]
    \tikzstyle{frontend}=[rectangle,rounded corners,draw=lightblue,fill=background,thick,minimum width=2.7cm,minimum height=1.3cm,text width=2.5cm,align=center,font=\small]
    \tikzstyle{backend}=[rectangle,rounded corners,draw=accent,fill=background,thick,minimum width=2.7cm,minimum height=1.3cm,text width=2.5cm,align=center,font=\small]
    \tikzstyle{db}=[cylinder,draw=darkblue,fill=background,thick,minimum width=1.6cm,minimum height=1.1cm,text width=1.4cm,align=center,font=\small]
    \tikzstyle{external}=[rectangle,rounded corners,draw=orange,fill=background,thick,minimum width=2.2cm,minimum height=1.1cm,text width=2cm,align=center,font=\small]
    \tikzstyle{arrow}=[->,>=stealth,line width=1.5pt]
    
    % Layers
    \node[client] (browser) {Web Browser\\ \small(HTTPS)};
    
    \node[frontend,below=2.5cm of browser] (frontend) {Angular\\ Frontend\\ :4200};
    
    \node[backend,below=2.5cm of frontend] (backend) {Express\\ Backend\\ :3000};
    
    \node[db,below left=2.5cm and 2.5cm of backend] (mongodb) {MongoDB\\ :27017};
    \node[db,below right=2.5cm and 2.5cm of backend] (camunda) {Camunda\\ :8080};
    
    \node[external,below left=2.5cm and 5cm of backend] (orangehrm) {OrangeHRM};
    \node[external,below=4cm of backend] (opencrx) {OpenCRX};
    \node[external,below right=2.5cm and 5cm of backend] (odoo) {Odoo};
    
    % Draw arrows
    \draw[arrow] (browser) -- node[right,text=darkblue] {JWT Token} (frontend);
    \draw[arrow] (frontend) -- node[right,text=lightblue] {REST API} (backend);
    \draw[arrow] (backend) -- (mongodb);
    \draw[arrow] (backend) -- (camunda);
    \draw[arrow] (backend) -- (orangehrm);
    \draw[arrow] (backend) -- (opencrx);
    \draw[arrow] (backend) -- (odoo);
    
\end{tikzpicture}
\caption{System Architecture Overview}
\label{fig:architecture}
\end{figure}

\section{Component Architecture}

\begin{figure}[H]
\centering
\begin{tikzpicture}[node distance=0.5cm and 0.7cm]
    
    % Controllers layer - draw box first as background
    \draw[rectangle,rounded corners,draw=darkblue,thick,fill=background] 
        (0,0) rectangle (9,2);
    \node[font=\small\bfseries,color=darkblue] at (4.5,1.7) {Controllers};
    
    % Components below the label
    \node[rectangle,rounded corners,draw=lightblue,fill=white,thick,minimum width=1.9cm,minimum height=0.7cm,font=\footnotesize] 
        at (1.2,0.8) {Auth};
    \node[rectangle,rounded corners,draw=lightblue,fill=white,thick,minimum width=1.9cm,minimum height=0.7cm,font=\footnotesize] 
        at (3.4,0.8) {Salesmen};
    \node[rectangle,rounded corners,draw=lightblue,fill=white,thick,minimum width=1.9cm,minimum height=0.7cm,font=\footnotesize] 
        at (5.6,0.8) {Bonus};
    \node[rectangle,rounded corners,draw=lightblue,fill=white,thick,minimum width=1.9cm,minimum height=0.7cm,font=\footnotesize] 
        at (7.8,0.8) {Performance};
    
    % Arrow from Controllers to Middleware
    \draw[->,thick] (4.5,0) -- (4.5,-0.3);
    
    % Middleware layer
    \draw[rectangle,rounded corners,draw=darkblue,thick,fill=background] 
        (0,-2.3) rectangle (9,-0.3);
    \node[font=\small\bfseries,color=darkblue] at (4.5,-0.6) {Middleware};
    
    \node[rectangle,rounded corners,draw=lightblue,fill=white,thick,minimum width=1.9cm,minimum height=0.7cm,font=\footnotesize] 
        at (2.3,-1.4) {JWT Auth};
    \node[rectangle,rounded corners,draw=lightblue,fill=white,thick,minimum width=1.9cm,minimum height=0.7cm,font=\footnotesize] 
        at (4.5,-1.4) {RBAC};
    \node[rectangle,rounded corners,draw=lightblue,fill=white,thick,minimum width=1.9cm,minimum height=0.7cm,font=\footnotesize] 
        at (6.7,-1.4) {Logging};
    
    % Arrow from Middleware to Services
    \draw[->,thick] (4.5,-2.3) -- (4.5,-2.6);
    
    % Services layer
    \draw[rectangle,rounded corners,draw=darkblue,thick,fill=background] 
        (0,-4.6) rectangle (9,-2.6);
    \node[font=\small\bfseries,color=darkblue] at (4.5,-2.9) {Services};
    
    \node[rectangle,rounded corners,draw=lightblue,fill=white,thick,minimum width=1.9cm,minimum height=0.7cm,font=\footnotesize] 
        at (1.2,-3.7) {Auth Svc};
    \node[rectangle,rounded corners,draw=lightblue,fill=white,thick,minimum width=1.9cm,minimum height=0.7cm,font=\footnotesize] 
        at (3.4,-3.7) {Bonus Svc};
    \node[rectangle,rounded corners,draw=lightblue,fill=white,thick,minimum width=1.9cm,minimum height=0.7cm,font=\footnotesize] 
        at (5.6,-3.7) {Cache Svc};
    \node[rectangle,rounded corners,draw=lightblue,fill=white,thick,minimum width=1.9cm,minimum height=0.7cm,font=\footnotesize] 
        at (7.8,-3.7) {Health Svc};
    
    % Arrow from Services to Integrations
    \draw[->,thick] (4.5,-4.6) -- (4.5,-4.9);
    
    % Integrations layer
    \draw[rectangle,rounded corners,draw=darkblue,thick,fill=background] 
        (0,-6.9) rectangle (9,-4.9);
    \node[font=\small\bfseries,color=darkblue] at (4.5,-5.2) {Integrations};
    
    \node[rectangle,rounded corners,draw=lightblue,fill=white,thick,minimum width=1.9cm,minimum height=0.7cm,font=\footnotesize] 
        at (1.2,-6) {OrangeHRM};
    \node[rectangle,rounded corners,draw=lightblue,fill=white,thick,minimum width=1.9cm,minimum height=0.7cm,font=\footnotesize] 
        at (3.4,-6) {OpenCRX};
    \node[rectangle,rounded corners,draw=lightblue,fill=white,thick,minimum width=1.9cm,minimum height=0.7cm,font=\footnotesize] 
        at (5.6,-6) {Odoo};
    \node[rectangle,rounded corners,draw=lightblue,fill=white,thick,minimum width=1.9cm,minimum height=0.7cm,font=\footnotesize] 
        at (7.8,-6) {Camunda};
    
\end{tikzpicture}
\caption{Express Backend Layer Architecture}
\label{fig:backend_layers}
\end{figure}

\section{Data Flow Diagram}

\begin{figure}[H]
\centering
\begin{tikzpicture}[node distance=3cm and 2cm]
    \tikzstyle{process}=[rectangle,rounded corners,draw=darkblue,fill=background,thick,minimum width=2.8cm,minimum height=1.2cm,text width=2.5cm,align=center,font=\small]
    \tikzstyle{data}=[ellipse,draw=lightblue,fill=background,thick,minimum width=2.5cm,minimum height=1.2cm,text width=2.2cm,align=center,font=\small]
    \tikzstyle{decision}=[diamond,draw=accent,fill=background,thick,minimum width=2cm,minimum height=2cm,text width=1.5cm,align=center,font=\small]
    \tikzstyle{arrow}=[->,>=stealth,line width=1.5pt]
    
    % User login flow
    \node[process] (login) {Login};
    \node[data,right=of login] (user_db) {User DB};
    \draw[arrow] (login) -- (user_db);
    
    % Bonus computation flow
    \node[process,below=of login] (compute) {Compute Bonus};
    \node[data,right=of compute] (perf_data) {Performance Data};
    \node[data,below=of perf_data] (order_data) {Order Data};
    \draw[arrow] (compute) -- (perf_data);
    \draw[arrow] (compute) -- (order_data);
    
    % Approval flow
    \node[decision,below=of compute] (approve) {Approve?};
    \node[process,below left=2cm and 1cm of approve] (ceo_approve) {CEO Approval};
    \node[process,below right=2cm and 1cm of approve] (hr_approve) {HR Approval};
    \draw[arrow] (compute) -- (approve);
    \draw[arrow] (approve) -- (ceo_approve);
    \draw[arrow] (approve) -- (hr_approve);
    
    % Storage
    \node[data,below=3cm of approve] (bonus_db) {Store Bonus};
    \draw[arrow] (ceo_approve) -- (bonus_db);
    \draw[arrow] (hr_approve) -- (bonus_db);
    
\end{tikzpicture}
\caption{Bonus Computation Data Flow}
\label{fig:dataflow}
\end{figure}

% ===== CHAPTER 3: PROJECT STRUCTURE =====
\chapter{Project Structure}

\section{Directory Hierarchy}

\begin{lstlisting}[language=bash,caption=Project Directory Layout,basicstyle=\ttfamily\footnotesize]
MaiPhucSan-InAr-develop/
  backend/                    # Node.js Express API
    src/
      controllers/            # HTTP request handlers
      middleware/             # Express middleware
      models/                 # Mongoose schemas
      routes/                 # API routes
      services/               # Business logic
      db/                     # Database config
      config/                 # Configuration
      app.js
    test/                     # Test files
    package.json
  frontend/                   # Angular application
    src/
      app/                    # Components & services
      environments/           # Config files
      main.ts
    angular.json
  docs/                       # Documentation
  workflow/                   # BPMN processes
  docker-compose.yml          # Infrastructure
\end{lstlisting}

% ===== CHAPTER 4: DATABASE MODELS =====
\chapter{Database Models}

\section{Data Model Overview}

\begin{figure}[H]
\centering
\begin{tikzpicture}[node distance=3.5cm and 3cm]
    \tikzstyle{entity}=[rectangle,draw=darkblue,fill=background,thick,minimum width=3.2cm,minimum height=2.8cm,text width=3cm,align=left,font=\small]
    \tikzstyle{relation}=[-,line width=1.5pt,color=accent]
    
    \node[entity] (user) {
        \textbf{\small User}\\[0.2cm]
        \rule{\linewidth}{0.5pt}\\[0.15cm]
        \fontsize{8}{9}\selectfont
        \_id: OID\\[0.05cm]
        username: Str\\[0.05cm]
        role: Enum\\[0.05cm]
        email: Str\\[0.05cm]
        hash: Str
    };
    
    \node[entity,right=of user] (salesman) {
        \textbf{\small Salesman}\\[0.2cm]
        \rule{\linewidth}{0.5pt}\\[0.15cm]
        \fontsize{8}{9}\selectfont
        \_id: OID\\[0.05cm]
        empId: Str\\[0.05cm]
        firstName: Str\\[0.05cm]
        lastName: Str\\[0.05cm]
        email: Str\\[0.05cm]
        manager: OID
    };
    
    \node[entity,below=of salesman] (performance) {
        \textbf{\small SocialPerf}\\[0.2cm]
        \rule{\linewidth}{0.5pt}\\[0.15cm]
        \fontsize{8}{9}\selectfont
        \_id: OID\\[0.05cm]
        empId: Str\\[0.05cm]
        month: Str\\[0.05cm]
        target: Num\\[0.05cm]
        actual: Num\\[0.05cm]
        bonus: Num
    };
    
    \node[entity,below=of user] (bonus) {
        \textbf{\small BonusComp}\\[0.2cm]
        \rule{\linewidth}{0.5pt}\\[0.15cm]
        \fontsize{8}{9}\selectfont
        \_id: OID\\[0.05cm]
        empId: Str\\[0.05cm]
        year: Num\\[0.05cm]
        totalBonus: Num\\[0.05cm]
        status: Enum\\[0.05cm]
        approvals: Obj
    };
    
    \node[entity,left=of bonus] (order) {
        \textbf{\small OrderEval}\\[0.2cm]
        \rule{\linewidth}{0.5pt}\\[0.15cm]
        \fontsize{8}{9}\selectfont
        \_id: OID\\[0.05cm]
        empId: Str\\[0.05cm]
        orderId: Str\\[0.05cm]
        revenueEur: Num\\[0.05cm]
        bonus: Num
    };
    
    % Relationships
    \draw[relation] (user) -- (salesman);
    \draw[relation] (salesman) -- (performance);
    \draw[relation] (salesman) -- (bonus);
    \draw[relation] (salesman) -- (order);
    
\end{tikzpicture}
\caption{Entity Relationship Diagram}
\label{fig:erd}
\end{figure}

\section{User Model}

\noindent\begin{minipage}{\linewidth}
\begin{lstlisting}[caption=User Schema Definition]
{
  _id: ObjectId,
  username: String (unique),
  email: String,
  passwordHash: String,
  role: Enum ["CEO", "HR", "SALESMAN", "ADMIN"],
  employeeId: String,
  firstName: String,
  lastName: String,
  isActive: Boolean,
  createdAt: Date,
  updatedAt: Date
}
\end{lstlisting}
\end{minipage}

\section{Bonus Computation Model}

\begin{lstlisting}[caption=BonusComputation Schema]
{
  _id: ObjectId,
  employeeId: String,
  year: Number,
  socialBonus: Number,
  orderBonus: Number,
  totalBonus: Number,
  status: Enum ["PENDING_CEO", "PENDING_HR", "RELEASED"],
  approvedByCEO: { date: Date, remarks: String },
  approvedByHR: { date: Date, remarks: String },
  storedInOrangeHRM: Boolean,
  workflowInstanceId: String (Camunda),
  createdAt: Date,
  updatedAt: Date
}
\end{lstlisting}

% ===== CHAPTER 5: API REFERENCE =====
\chapter{API Reference}

\section{Authentication}

\subsection{Login Endpoint}

\endpoint{POST /api/v1/auth/login}

\begin{lstlisting}[caption=Login Request/Response]
# Request
POST /api/v1/auth/login
Content-Type: application/json

{
  "username": "ceo",
  "password": "Ceo123!"
}

# Response (200)
{
  "token": "eyJhbGciOiJIUzI1NiIs...",
  "user": {
    "id": "user_123",
    "username": "ceo",
    "role": "CEO",
    "employeeId": "E001"
  }
}
\end{lstlisting}

\section{Salesmen Management}

\subsection{List Salesmen}

\endpoint{GET /api/v1/salesmen}

\begin{table}[H]
\centering
\begin{tabular}{lll}
\toprule
\textbf{Parameter} & \textbf{Type} & \textbf{Description} \\
\midrule
Authorization & Header & Bearer token (required) \\
page & Query & Page number (default: 1) \\
limit & Query & Results per page (default: 20) \\
\bottomrule
\end{tabular}
\caption{Salesmen List Parameters}
\end{table}

\subsection{Create Salesman}

\endpoint{POST /api/v1/salesmen}

\begin{lstlisting}[caption=Create Salesman Request]
POST /api/v1/salesmen
Authorization: Bearer <token>
Content-Type: application/json

{
  "firstName": "John",
  "lastName": "Doe",
  "employeeId": "E1001",
  "email": "john@company.com",
  "department": "Sales"
}
\end{lstlisting}

\section{Bonus Computation}

\subsection{Compute Bonus}

\endpoint{POST /api/v1/bonus/:employeeId/compute}

This endpoint calculates the bonus based on performance records and orders.

\begin{lstlisting}[caption=Bonus Computation Request]
POST /api/v1/bonus/E1001/compute
Authorization: Bearer <token>
Content-Type: application/json

{
  "year": 2025,
  "reason": "Q1 Final Calculation"
}

# Response (200)
{
  "employeeId": "E1001",
  "year": 2025,
  "socialBonus": 18000,
  "orderBonus": 5000,
  "totalBonus": 23000,
  "status": "PENDING_CEO"
}
\end{lstlisting}

\subsection{Approval Endpoints}

\begin{table}[H]
\centering
\begin{tabular}{ll}
\toprule
\textbf{Endpoint} & \textbf{Purpose} \\
\midrule
\code{POST /api/v1/bonus/:employeeId/approve/ceo} & CEO approval \\
\code{POST /api/v1/bonus/:employeeId/approve/hr} & HR approval \& storage \\
\code{GET /api/v1/bonus/:employeeId} & Get bonus details \\
\code{GET /api/v1/bonus/:employeeId/history} & View history \\
\bottomrule
\end{tabular}
\caption{Bonus API Endpoints}
\end{table}

% ===== CHAPTER 6: WORKFLOW & INTEGRATION =====
\chapter{Workflow Orchestration}

\section{Bonus Approval Workflow}

\begin{figure}[H]
\centering
\begin{tikzpicture}[node distance=1.8cm and 1.5cm]
    \tikzstyle{state}=[rectangle,rounded corners,draw=darkblue,fill=background,thick,minimum width=2.5cm,minimum height=1cm,text width=2.3cm,align=center,font=\footnotesize]
    \tikzstyle{decision}=[diamond,draw=accent,fill=background,thick,minimum width=1.8cm,minimum height=1.8cm,text width=1.5cm,align=center,font=\footnotesize]
    \tikzstyle{arrow}=[->,>=stealth,line width=1.2pt,color=darkblue]
    
    \node[state] (init) {Initialize};
    \node[state,above=of init] (ceo) {CEO Review};
    \node[decision,above=of ceo] (ceo_dec) {Approved?};
    \node[state,right=2.5cm of ceo_dec] (hr) {HR Review};
    \node[decision,above=of hr] (hr_dec) {Approved?};
    \node[state,above=of hr_dec] (store) {Store in OrangeHRM};
    \node[state,above=of store] (release) {Release to Salesman};
    \node[state,left=2.5cm of ceo_dec] (reject) {Rejected};
    
    \draw[arrow] (init) -- (ceo);
    \draw[arrow] (ceo) -- (ceo_dec);
    \draw[arrow] (ceo_dec) -- node[above,font=\tiny] {✓ Yes} (hr);
    \draw[arrow] (ceo_dec) -- node[above,font=\tiny] {✗ No} (reject);
    \draw[arrow] (hr) -- (hr_dec);
    \draw[arrow] (hr_dec) -- node[right,font=\tiny] {✓ Yes} (store);
    \draw[arrow] (hr_dec.west) -- node[above,font=\tiny] {✗ No} (reject.north);
    \draw[arrow] (store) -- (release);
    
\end{tikzpicture}
\caption{Bonus Approval BPMN Workflow}
\label{fig:workflow}
\end{figure}

\section{Camunda Integration}

The system uses Camunda BPM for process orchestration. Key features:

\begin{itemize}
    \item Workflow instantiation via \code{CamundaService}
    \item Process variable management
    \item Task assignment to roles
    \item Audit trail for all process steps
\end{itemize}

\note{The BPMN file is located at \code{workflow/bonusApproval.bpmn} and can be deployed via Camunda Modeler or REST API.}

% ===== CHAPTER 7: SETUP & DEPLOYMENT =====
\chapter{Setup \& Deployment}

\section{Prerequisites}

\begin{tcolorbox}[colback=background,colframe=darkblue,title=System Requirements]
\begin{itemize}
    \item Node.js 18+ (18.18.0 recommended)
    \item npm 9+
    \item Docker \& Docker Compose
    \item Git
    \item 2GB RAM minimum
    \item 5GB disk space minimum
\end{itemize}
\end{tcolorbox}

\section{Development Setup}

\subsection{Step 1: Infrastructure}

\begin{lstlisting}[language=bash,caption=Start MongoDB and Camunda]
# From project root
docker-compose up -d

# Verify services
docker-compose ps
\end{lstlisting}

\subsection{Step 2: Backend Setup}

\begin{lstlisting}[language=bash,caption=Backend Installation]
cd backend
npm install
cp ../.env.example .env
npm run seed           # Seed demo data
npm run dev            # Start dev server
# API available at http://localhost:3000/api/v1
\end{lstlisting}

\subsection{Step 3: Frontend Setup}

\begin{lstlisting}[language=bash,caption=Frontend Installation]
cd frontend
npm install
npm start              # Start dev server
# Frontend available at http://localhost:4200
\end{lstlisting}

\section{Production Deployment}

\subsection{Docker Build}

\begin{lstlisting}[language=dockerfile,caption=Backend Dockerfile]
FROM node:18-alpine

WORKDIR /app

COPY package*.json ./
RUN npm ci --only=production

COPY src ./src
COPY openapi ./openapi

EXPOSE 3000
CMD ["node", "src/server.js"]
\end{lstlisting}

\subsection{Environment Configuration}

\begin{table}[H]
\centering
\small
\begin{tabular}{|l|l|p{4cm}|}
\hline
\textbf{Variable} & \textbf{Example} & \textbf{Description} \\
\hline
NODE\_ENV & production & Environment mode \\
PORT & 3000 & Server port \\
MONGODB\_URI & mongodb://localhost:27017 & Database URL \\
JWT\_SECRET & your-secret-key & JWT signing secret \\
CORS\_ORIGIN & https://app.example.com & CORS allowed origin \\
CAMUNDA\_URL & http://camunda:8080 & Camunda engine URL \\
\hline
\end{tabular}
\caption{Environment Variables}
\end{table}

% ===== CHAPTER 8: INTEGRATIONS =====
\chapter{External Integrations}

\section{OrangeHRM Integration}

\begin{figure}[H]
\centering
\begin{tikzpicture}[node distance=4cm and 2cm]
    \tikzstyle{box}=[rectangle,rounded corners,draw=darkblue,fill=background,thick,minimum width=3cm,minimum height=1.2cm,text width=2.8cm,align=center,font=\small]
    \tikzstyle{arrow}=[->,>=stealth,line width=1.5pt]
    
    \node[box] (app) {HighPerformance};
    \node[box,right=of app] (orangehrm) {OrangeHRM};
    
    \draw[arrow] (app) -- node[above,font=\small] {REST API} (orangehrm);
    \draw[arrow] (orangehrm) -- node[below,font=\small] {JSON} (app);
    
    \node[below=1.5cm of app,align=left,font=\small] (ops) {
        • Fetch Employee Data \\
        • Store Bonuses \\
        • Update Custom Fields
    };
    
\end{tikzpicture}
\caption{OrangeHRM Integration Flow}
\end{figure}

\subsection{Configuration}

\begin{lstlisting}[caption=OrangeHRM Setup]
# .env
ORANGEHRM_BASE_URL=https://orangehrm.example.com
ORANGEHRM_USERNAME=admin
ORANGEHRM_PASSWORD=password
\end{lstlisting}

\section{OpenCRX Integration}

Orders are fetched from OpenCRX and cached for bonus calculation. The integration provides:

\begin{itemize}
    \item Order data retrieval (JSON-RPC)
    \item Client ranking information
    \item Revenue and probability metrics
    \item Bonus calculation from orders
\end{itemize}

\section{Odoo Integration}

\noindent\begin{minipage}{\linewidth}
\begin{lstlisting}[caption=Odoo JSON-RPC Integration]
# Fetch employee data
POST /jsonrpc
{
  "jsonrpc": "2.0",
  "method": "call",
  "params": {
    "service": "object",
    "method": "execute",
    "args": ["database", "uid", "password", 
             "hr.employee", "search_read", [], 
             ["name", "email", "department_id"]]
  }
}
\end{lstlisting}
\end{minipage}

% ===== CHAPTER 9: TESTING =====
\chapter{Testing \& Quality}

\section{Test Architecture}

\begin{figure}[H]
\centering
\begin{tikzpicture}[node distance=3cm and 2cm]
    \tikzstyle{test}=[rectangle,rounded corners,draw=darkblue,fill=background,thick,minimum width=2.8cm,minimum height=1.2cm,text width=2.5cm,align=center,font=\small]
    
    \node[test] (unit) {Unit Tests};
    \node[test,right=of unit] (integration) {Integration Tests};
    \node[test,right=of integration] (api) {API Tests};
    
    \node[below=of unit,font=\small] (mocha) {Mocha};
    \node[below=of integration,font=\small] (chai) {Chai + Sinon};
    \node[below=of api,font=\small] (postman) {Postman};
    
    \draw[->,thick] (unit) -- (mocha);
    \draw[->,thick] (integration) -- (chai);
    \draw[->,thick] (api) -- (postman);
    
\end{tikzpicture}
\caption{Testing Framework Stack}
\end{figure}

\section{Running Tests}

\begin{lstlisting}[language=bash,caption=Test Execution]
# Unit and integration tests
cd backend
npm test

# Coverage report
npm test -- --coverage

# Watch mode
npm test -- --watch
\end{lstlisting}

\section{Test Suites}

\begin{table}[H]
\centering
\begin{tabular}{lll}
\toprule
\textbf{Suite} & \textbf{File} & \textbf{Coverage} \\
\midrule
Bonus Service & \code{bonus.service.test.js} & Calculations \\
Dependencies & \code{dependencies.test.js} & Integrations \\
Integration & \code{integration/} & Full workflows \\
\bottomrule
\end{tabular}
\caption{Test Suite Organization}
\end{table}

% ===== CHAPTER 10: TROUBLESHOOTING =====
\chapter{Troubleshooting \& Support}

\section{Common Issues}

\subsection{MongoDB Connection Failed}

\begin{tcolorbox}[colback=orange!5!white,colframe=orange,title=Issue: Cannot Connect to MongoDB]
\textbf{Solution:}
\begin{enumerate}
    \item Check if MongoDB container is running: \code{docker-compose ps}
    \item Restart MongoDB: \code{docker-compose restart mongodb}
    \item Verify URI in .env file
    \item Check port 27017 availability
\end{enumerate}
\end{tcolorbox}

\subsection{Backend Won't Start}

\begin{lstlisting}[language=bash,caption=Debugging Backend Issues]
# Check port conflicts
netstat -ano | findstr :3000

# Clear dependencies and reinstall
rm -rf node_modules package-lock.json
npm install
npm run dev

# Check logs
npm run dev 2>&1 | head -100
\end{lstlisting}

\subsection{Frontend Build Errors}

\begin{lstlisting}[language=bash,caption=Fixing Angular Build Issues]
# Clear Angular cache
rm -rf node_modules/.angular

# Rebuild
npm install
npm start
\end{lstlisting}

\section{Performance Optimization}

\begin{itemize}
    \item \textbf{Database Indexing} --- Create indexes on frequently queried fields
    \item \textbf{Response Compression} --- Enable gzip in Express
    \item \textbf{Caching} --- Implement Redis for session management
    \item \textbf{API Rate Limiting} --- Prevent abuse
    \item \textbf{Connection Pooling} --- MongoDB connection optimization
\end{itemize}

\section{Security Checklist}

\begin{tcolorbox}[colback=background,colframe=darkblue,title=Production Security Checklist]
\begin{itemize}
    \item[$\square$] Change JWT\_SECRET to strong random value
    \item[$\square$] Use HTTPS/TLS in production
    \item[$\square$] Secure MongoDB with authentication
    \item[$\square$] Implement rate limiting
    \item[$\square$] Enable CORS properly
    \item[$\square$] Validate all user inputs
    \item[$\square$] Use environment variables for secrets
    \item[$\square$] Regular dependency updates
    \item[$\square$] Enable audit logging
    \item[$\square$] Rotate credentials periodically
\end{itemize}
\end{tcolorbox}

% ===== CHAPTER 11: BONUS CALCULATION =====
\chapter{Bonus Calculation Algorithm}

\section{Overview}

The bonus calculation system computes bonuses based on two components:

\begin{enumerate}
    \item \textbf{Social Performance Bonus} --- Based on actual vs. target metrics
    \item \textbf{Order-Based Bonus} --- Based on closed deals and revenue
\end{enumerate}

\section{Social Performance Calculation}

\subsection{Formula}

The social performance bonus is calculated using:

\[
\text{Bonus}_{\text{social}} = \sum_{i=1}^{n} \frac{(\text{Actual}_i - \text{Target}_i)}{\text{Target}_i} \times 100,000
\]

Where:
\begin{itemize}
    \item $\text{Actual}_i$ = Actual achieved metric for period $i$
    \item $\text{Target}_i$ = Target metric for period $i$
    \item $100,000$ = Base bonus amount
\end{itemize}

\subsection{Example}

\begin{table}[H]
\centering
\begin{tabular}{|c|c|c|c|}
\hline
\textbf{Month} & \textbf{Target} & \textbf{Actual} & \textbf{Bonus Contribution} \\
\hline
January & 100,000 & 120,000 & 20,000 EUR \\
February & 100,000 & 110,000 & 10,000 EUR \\
March & 100,000 & 100,000 & 0 EUR \\
\hline
Total & 300,000 & 330,000 & 30,000 EUR \\
\hline
\end{tabular}
\caption{Social Performance Bonus Example}
\end{table}

\subsection{Implementation}

\noindent\begin{minipage}{\linewidth}
\begin{lstlisting}[caption=Bonus Calculation Implementation]
async function computeSocialRecordBonus(record) {
  const { target, actual } = record;
  if (!target || target === 0) return 0;
  
  const variance = (actual - target) / target;
  const bonus = variance * 100000;
  
  return Math.max(0, Math.round(bonus));
}
\end{lstlisting}
\end{minipage}

\section{Order-Based Bonus Calculation}

\subsection{Formula}

Order bonuses consider revenue and closing probability:

\[
\text{Bonus}_{\text{order}} = \text{Revenue} \times \text{Probability} \times 0.05
\]

Where:
\begin{itemize}
    \item \text{Revenue} = Order value in EUR
    \item \text{Probability} = Closing probability (0-100\%)
    \item $0.05$ = Commission rate
\end{itemize}

\subsection{Example}

\begin{table}[H]
\centering
\begin{tabular}{|c|c|c|c|}
\hline
\textbf{Order} & \textbf{Revenue} & \textbf{Probability} & \textbf{Bonus} \\
\hline
ORD001 & 50,000 EUR & 95\% & 2,375 EUR \\
ORD002 & 30,000 EUR & 80\% & 1,200 EUR \\
ORD003 & 20,000 EUR & 60\% & 600 EUR \\
\hline
Total & 100,000 EUR & --- & 4,175 EUR \\
\hline
\end{tabular}
\caption{Order Bonus Calculation Example}
\end{table}

\section{Total Bonus Computation}

\[
\text{Total Bonus} = \text{Bonus}_{\text{social}} + \text{Bonus}_{\text{order}}
\]

The system persists this as a \code{BonusComputation} record with status tracking through the approval workflow.

\section{Bonus Constraints}

\begin{itemize}
    \item Minimum bonus: 0 EUR (no negative bonuses)
    \item Maximum bonus: 500,000 EUR (configurable)
    \item Precision: 2 decimal places
    \item Year-based: Annual calculations with monthly/quarterly breakdown
\end{itemize}

% ===== CHAPTER 12: WORKFLOW & APPROVAL PROCESS =====
\chapter{Workflow \& Approval Process}

\section{BPMN Workflow Overview}

The bonus approval workflow is implemented using Camunda BPM with the following states:

\subsection{Workflow State Machine}

\begin{figure}[H]
\centering
\begin{tikzpicture}[
    node distance=1.8cm and 1.5cm,
    state/.style={circle, draw, minimum size=1.5cm, text centered, font=\footnotesize},
    decision/.style={diamond, draw, aspect=2, text centered, minimum width=1.8cm, font=\footnotesize},
    terminal/.style={rounded rectangle, draw, minimum width=1.8cm, minimum height=0.7cm, font=\footnotesize}
]

% Main flow - vertical
\node[state] (draft) {Draft};
\node[state, above=of draft] (pending) {Pending};
\node[decision, above=of pending] (ceo) {CEO\\Review};
\node[state, above=of ceo] (ceoapproved) {CEO\\Approved};
\node[decision, above=of ceoapproved] (hr) {HR\\Review};

% Terminals - side by side at top
\node[terminal, left=2cm of hr] (approved) {Approved};
\node[terminal, right=2cm of hr] (rejected) {Rejected};

% Main flow arrows
\draw[->, thick] (draft) -- node[right,font=\tiny] {Submit} (pending);
\draw[->, thick] (pending) -- node[right,font=\tiny] {Send} (ceo);
\draw[->, thick] (ceo) -- node[right,font=\tiny] {Approve} (ceoapproved);
\draw[->, thick] (ceoapproved) -- node[right,font=\tiny] {Forward} (hr);
\draw[->, thick] (hr) -- node[above,font=\tiny] {Approve} (approved);
\draw[->, thick] (hr) -- node[above,font=\tiny] {Reject} (rejected);

% Resubmit loop - go around the side
\draw[->, thick] (rejected) -- ++(1.5,0) |- node[near start,right,font=\tiny] {Resubmit} (draft);

\end{tikzpicture}
\caption{Bonus Approval Workflow}
\end{figure}

\section{Workflow States}

\begin{description}[labelindent=1cm, font=\textbf]
    \item[Draft] - Initial state when bonus computation is created
    \item[Pending] - Submitted for CEO review
    \item[CEO Approved] - CEO has reviewed and approved
    \item[HR Approved] - Final approval from HR (Approved state)
    \item[Rejected] - Rejected by CEO or HR, returned to Draft
\end{description}

\section{Approval Task Details}

\subsection{CEO Review Task}

\begin{tcolorbox}[title=CEO Review Process]
\begin{itemize}
    \item \textbf{Owner:} CEO (role-based assignment)
    \item \textbf{Deadline:} 5 business days
    \item \textbf{Actions:}
    \begin{itemize}
        \item Review bonus computation summary
        \item Check total bonus amount against budget
        \item Approve or reject with comments
    \end{itemize}
    \item \textbf{Exit Criteria:}
    \begin{itemize}
        \item Approved → forward to HR
        \item Rejected → return to draft with feedback
    \end{itemize}
\end{itemize}
\end{tcolorbox}

\subsection{HR Review Task}

\begin{tcolorbox}[title=HR Final Approval Process]
\begin{itemize}
    \item \textbf{Owner:} HR Team (role-based assignment)
    \item \textbf{Deadline:} 3 business days
    \item \textbf{Actions:}
    \begin{itemize}
        \item Verify all data completeness
        \item Check compliance with company policies
        \item Final approval or rejection
    \end{itemize}
    \item \textbf{Exit Criteria:}
    \begin{itemize}
        \item Approved → bonus marked as finalized
        \item Rejected → return to draft
    \end{itemize}
\end{itemize}
\end{tcolorbox}

\section{Workflow API Endpoints}

\begin{table}[H]
\centering
\small
\begin{tabular}{|l|l|l|p{3cm}|}
\hline
\textbf{Method} & \textbf{Endpoint} & \textbf{Action} & \textbf{Description} \\
\hline
POST & \texttt{/api/workflow/\{id\}/approve} & CEO/HR Approve & Approve bonus computation \\
POST & \texttt{/api/workflow/\{id\}/reject} & Reject & Return to draft with reason \\
GET & \texttt{/api/workflow/\{id\}/history} & Get History & View state transitions \\
GET & \texttt{/api/workflow/tasks} & Get Tasks & List pending tasks \\
\hline
\end{tabular}
\caption{Workflow API Endpoints}
\end{table}

\section{Integration with Services}

\subsection{Camunda Service Integration}

The workflow is managed through the \code{CamundaService}:

\begin{lstlisting}[caption=Camunda Service Integration Example]
// Start workflow
await camundaService.startProcess('bonus-approval', {
  bonusId: computation.id,
  totalAmount: computation.total,
  salesman: computation.salesman
});

// Complete CEO approval task
await camundaService.completeTask(taskId, {
  action: 'approve',
  comments: 'Approved by CEO',
  userId: req.user.id
});
\end{lstlisting}

\section{Notifications \& Escalations}

\subsection{Email Notifications}

\begin{itemize}
    \item \textbf{To CEO:} When bonus reaches pending state
    \item \textbf{To HR:} When CEO approves
    \item \textbf{To Salesman:} Final approval notification
    \item \textbf{Escalation:} If approval exceeds deadline
\end{itemize}

\subsection{Escalation Rules}

\begin{tcolorbox}[colframe=red!70, title=Escalation Triggers]
\begin{itemize}
    \item CEO review exceeds 5 days → escalate to CFO
    \item HR review exceeds 3 days → escalate to Head of HR
    \item Rejected twice → manual review required
\end{itemize}
\end{tcolorbox}

\section{Workflow Monitoring}

Access workflow metrics at: \url{http://localhost:8080/camunda}

\begin{itemize}
    \item Active process instances
    \item Running tasks
    \item Completed milestones
    \item Performance metrics
\end{itemize}

% ===== APPENDIX =====
\appendix

\chapter{Quick Reference}

\section{API Endpoints Summary}

\begin{longtable}{lll}
\toprule
\textbf{Method} & \textbf{Endpoint} & \textbf{Purpose} \\
\midrule
\endfirsthead
\multicolumn{3}{c}{\textit{continued from previous page}} \\
\toprule
\textbf{Method} & \textbf{Endpoint} & \textbf{Purpose} \\
\midrule
\endhead
\endfoot
\hline
\endlastfoot
POST & /api/v1/auth/login & User login \\
GET & /api/v1/salesmen & List salesmen \\
POST & /api/v1/salesmen & Create salesman \\
GET & /api/v1/bonus/:id & Get bonus \\
POST & /api/v1/bonus/:id/compute & Compute bonus \\
POST & /api/v1/bonus/:id/approve/ceo & CEO approval \\
POST & /api/v1/bonus/:id/approve/hr & HR approval \\
GET & /api/v1/performance/social/:id & Get performance \\
POST & /api/v1/performance/social & Create record \\
GET & /api/v1/orders/:id & Get orders \\
GET & /api/v1/health-check & System health \\
\caption{API Endpoints Reference}
\label{tab:endpoints}
\end{longtable}

\section{Demo Credentials}

\begin{table}[H]
\centering
\begin{tabular}{lll}
\toprule
\textbf{Role} & \textbf{Username} & \textbf{Password} \\
\midrule
CEO & ceo & Ceo123! \\
HR & hr & Hr123! \\
Salesman & salesman01 & password123 \\
\bottomrule
\end{tabular}
\caption{Demo User Credentials}
\end{table}

\section{Useful Commands}

\noindent\begin{minipage}{\linewidth}
\begin{lstlisting}[language=bash,caption=Quick Reference Commands]
# Start infrastructure
docker-compose up -d

# Backend operations
cd backend && npm install && npm run seed && npm run dev

# Frontend operations
cd frontend && npm install && npm start

# Run tests
npm run backend:test

# Check health
curl http://localhost:3000/api/v1/health-check

# View API docs
# Open: http://localhost:3000/api-docs
\end{lstlisting}
\end{minipage}

\chapter{Glossary}

\begin{description}
    \item[BPMN] Business Process Model and Notation --- Standard for workflow design
    \item[RBAC] Role-Based Access Control --- Security model for authorization
    \item[JWT] JSON Web Token --- Stateless authentication method
    \item[REST] Representational State Transfer --- API architectural style
    \item[ORM] Object-Relational Mapping --- Mongoose provides this for MongoDB
    \item[JSON-RPC] JSON Remote Procedure Call --- Lightweight RPC protocol
    \item[OpenAPI] API specification standard for REST services
    \item[CORS] Cross-Origin Resource Sharing --- HTTP-based access control
\end{description}

% Index
\printindex

\end{document}
